% ============================================================
%  chapters/06_results.tex
% ============================================================

\section{Results}
\label{sec:results}

\subsection{School Enrollment}
\label{sec:results_attend}

Table~\ref{tab:attend_results} presents OLS estimates of
Equation~\eqref{eq:attend} (Column 1) and Equation~\eqref{eq:attend_fe}
with LGA fixed effects (Column 2) for the \textit{attend} outcome.

Looking at the baseline OLS results (Column 1), the violence indicators
statistically significant at conventional levels are exposure during early
childhood (ages 0--3), preschool age (3--6), and primary school age
(6--12). The coefficient on \textit{viol\_ech} is $-0.022$ ($p < 0.01$),
on \textit{viol\_ps} is $-0.049$ ($p < 0.001$), and on \textit{viol\_prim}
is $+0.019$ ($p < 0.01$). These estimates indicate that children exposed to
violence at younger ages are less likely to be currently enrolled in school,
while children who experienced violence during primary school age exhibit a
counterintuitive positive association with current enrollment. The remaining
violence variables carry positive coefficients except for
\textit{viol\_SS}, which is negative, though none are statistically
significant. Pre-birth violence exposure shows no statistically significant
effect on current enrollment.

When LGA fixed effects are included (Column 2), the model fit improves
substantially (R-squared increases from 0.43 to 0.52), and several
previously insignificant violence variables become significant. Exposure
during the pre-birth window (17--9 months before birth) is now positive
and significant at the 5\% level ($\hat{\beta} = 0.031$), and in-utero
exposure is positive and significant at the 1\% level ($\hat{\beta} =
0.039$). The early childhood effect is no longer statistically significant
and the coefficient falls to near zero. The primary school age effect
remains positive and significant, with a larger coefficient ($0.055$) than
in the baseline model.

Across both specifications, male students are significantly more likely to
be enrolled in school (coefficient $\approx 0.04$--$0.05$, $p < 0.001$).
The urban indicator is small and insignificant in the baseline but becomes
larger and significant with fixed effects ($\hat{\delta} = 0.035$), as
expected once we absorb time-invariant LGA-level confounders. All parental
education variables are positive, and mother's and father's alive status
are not statistically significant with LGA fixed effects. Distance to
primary school has a statistically significant negative effect on
enrollment; distance to junior secondary and senior secondary schools is
not significant.

% ---- Main regression table placeholder ---------------------
% Export this from Python (e.g., using stargazer or modelsummary)
% or Stata and paste here, or write it manually.

\begin{table}[H]
    \centering
    \caption{OLS Estimates: Effect of Violence Exposure on School Enrollment}
    \label{tab:attend_results}
    \begin{threeparttable}
        \begin{tabular}{lcc}
            \toprule
            & \multicolumn{2}{c}{Dependent variable: \textit{attend}} \\
            \cmidrule(lr){2-3}
            & (1) OLS & (2) OLS + LGA FE \\
            \midrule
            \multicolumn{3}{l}{\textit{Violence exposure indicators}} \\[0.3em]
            Pre-birth ($-$17 to $-$9 months) & 0.004 & 0.031$^{**}$ \\
                                          & (0.014) & (0.015) \\[0.3em]
            In utero (0–9 months pre-birth) & 0.016 & 0.039$^{***}$ \\
                                             & (0.013) & (0.014) \\[0.3em]
            Early childhood (0–3 years) & $-$0.022$^{***}$ & 0.010 \\
                                         & (0.007) & (0.010) \\[0.3em]
            Preschool age (3–6 years) & $-$0.049$^{***}$ & $-$0.017 \\
                                       & (0.007) & (0.010) \\[0.3em]
            Primary school (6–12 years) & 0.019$^{***}$ & 0.055$^{***}$ \\
                                         & (0.007) & (0.011) \\[0.3em]
            Junior secondary (12–15 years) & 0.018 & 0.032 \\
                                            & (0.019) & (0.023) \\[0.3em]
            Senior secondary (15–18 years) & $-$0.019 & $-$0.115 \\
                                            & (0.074) & (0.105) \\[0.5em]
            Individual controls & Yes & Yes \\
            LGA fixed effects   & No  & Yes \\
            \midrule
            Observations & 56,878 & 56,878 \\
            R-squared    & 0.428  & 0.517  \\
            \bottomrule
        \end{tabular}
        \begin{tablenotes}
            \small
            \item \textit{Notes:} OLS estimates with robust standard errors
            in parentheses. Individual controls include age, age-squared,
            gender, urban/rural, region, wealth quintile, parental
            education and living status, Quranic/Tsangaya/Islamiyya school
            attendance, late school start, disability indicators, and
            distances to school. $^{*}p < 0.10$, $^{**}p < 0.05$,
            $^{***}p < 0.01$.
        \end{tablenotes}
    \end{threeparttable}
\end{table}

\subsection{Educational Attainment}
\label{sec:results_attainment}

Table~\ref{tab:attainment_results} presents results for the
\textit{highest\_enrolled} outcome---total years of schooling completed.

In the baseline OLS (Column 1), nearly all violence regressors carry
negative coefficients, with the exceptions of \textit{viol\_ps} (preschool,
$\hat{\beta} = 0.154$, $p < 0.01$) and \textit{viol\_JS} (junior secondary,
$\hat{\beta} = 0.052$, not significant). All violence variables except
\textit{viol\_prim}, \textit{viol\_JS}, and \textit{viol\_SS} are
statistically significant.

The LGA fixed effects results (Column 2) tell a markedly different story.
After absorbing LGA-level confounders, all violence indicators except
\textit{viol\_SS} exhibit positive associations with educational attainment,
and most are statistically significant at the 5\% level or better
(\textit{viol\_uter\_p} and \textit{viol\_ech} are not significant). This
reversal in sign across specifications suggests substantial negative
selection of violence into low-performing LGAs: once we compare children
within the same LGA, violence exposure is associated with more---not
fewer---years of schooling. This result is difficult to interpret causally
and points to important identification challenges discussed in
Section~\ref{sec:conclusion}.

Male students achieve more years of schooling than female students across
both specifications (coefficient $\approx 0.014$, significant). Urban
children achieve significantly more schooling, and the urban premium triples
in magnitude with the introduction of LGA fixed effects. Similarly, the
coefficient on \textit{region} increases from approximately 0.06 to 0.21
with fixed effects, underscoring the importance of geographic controls.

% ---- Attainment regression table placeholder ----------------
\begin{table}[H]
    \centering
    \caption{OLS Estimates: Effect of Violence Exposure on Years of Schooling}
    \label{tab:attainment_results}
    \begin{threeparttable}
        \begin{tabular}{lcc}
            \toprule
            & \multicolumn{2}{c}{Dependent variable: \textit{highest\_enrolled}} \\
            \cmidrule(lr){2-3}
            & (1) OLS & (2) OLS + LGA FE \\
            \midrule
            \multicolumn{3}{l}{\textit{Violence exposure indicators}} \\[0.3em]
            Pre-birth ($-$17 to $-$9 months) & & \\[0.3em]
            In utero (0–9 months pre-birth) & & \\[0.3em]
            Early childhood (0–3 years) & & \\[0.3em]
            Preschool age (3–6 years) & $0.154^{***}$ & \\[0.3em]
            Primary school (6–12 years) & & \\[0.3em]
            Junior secondary (12–15 years) & $0.052$ & \\[0.3em]
            Senior secondary (15–18 years) & & \\[0.5em]
            Individual controls & Yes & Yes \\
            LGA fixed effects   & No  & Yes \\
            \midrule
            Observations & [N] & [N] \\
            R-squared    & & \\
            \bottomrule
        \end{tabular}
        \begin{tablenotes}
            \small
            \item \textit{Notes:} OLS estimates with robust standard errors
            in parentheses. See Table~\ref{tab:attend_results} notes for
            description of individual controls. Full regression output is
            in Appendix~\ref{app:full_results}.
            $^{*}p < 0.10$, $^{**}p < 0.05$, $^{***}p < 0.01$.
        \end{tablenotes}
    \end{threeparttable}
\end{table}

% ---- Placeholder: add figures here when ready ---------------
% Example figure call — see main.tex for instructions:
%
% \begin{figure}[H]
%     \centering
%     \includegraphics[width=0.85\textwidth]{figures/coef_plot_attend.pdf}
%     \caption{Violence Exposure Coefficients: School Enrollment}
%     \label{fig:coef_attend}
% \end{figure}
